\documentclass[11pt,letterpaper]{article}

\usepackage[margin=0.72in]{geometry}
\usepackage[T1]{fontenc}
\usepackage{lmodern}
\usepackage[hidelinks]{hyperref}
\usepackage{microtype}
\setlength{\parindent}{0pt}
\setlength{\parskip}{8pt}
\pagestyle{empty}

\begin{document}

\textbf{\Large Piter Z. Garcia Bautista}\\
Rochester, NY \enspace|\enspace
\href{mailto:pzg8794@rit.edu}{pzg8794@rit.edu} \enspace|\enspace
+1 (646) 288-3300\\
\href{https://linkedin.com/in/garciapiterz}{linkedin.com/in/garciapiterz}
\enspace|\enspace
\href{https://github.com/pzg8794}{github.com/pzg8794}

\vspace{10pt}
August 18, 2026

Hiring Team\\
DeepBrain Technology

\textbf{Re: Research Scientist / Research Fellow (AI for Education \& Game Intelligence)}

Dear DeepBrain Technology Hiring Team:

I am applying because this role connects the two areas I am deliberately
bringing together: applied AI research and inclusive K--12 computer science
education. I do not yet hold a Ph.D.; I am completing M.S. degrees in Data
Science at Rochester Institute of Technology and Teaching Computer Science
K--12 at the University of Rochester, both expected in December 2026, and I am
preparing to begin doctoral study in Fall 2027. I offer more than ten years of
professional engineering and data-systems experience, a completed M.S. in
Computer Science, current AI research, and classroom-based teacher preparation
as an experience-equivalency case for your consideration.

As a graduate assistant in RIT Software Engineering, I build reproducible
Python evaluation workflows for adaptive decision systems in quantum-network
routing. My work includes multi-armed, contextual, adversarial, and neural
bandit methods; automated validation; structured experiment logging; benchmark
analysis; and the translation of results into technical reports,
presentations, and manuscript-ready artifacts. Earlier engineering roles
developed my ability to build dependable data workflows, integrate machine
learning into production systems, and collaborate across technical and
nontechnical teams.

My education work gives that research a practical learning-sciences context.
Through K--5 computer science teaching, middle-school science teaching, an
environmental-science camp, and graduate study in inclusive and digitally rich
teaching, I design standards-aligned learning with Universal Design for
Learning, multiple participation pathways, and critical AI literacy. I am
especially interested in research on adaptive feedback, student modeling,
reasoning assessment, and the responsible use of AI as a learning partner.

I recognize that the posting requests a completed Ph.D. and a strong
publication record. I would welcome consideration if DeepBrain is open to an
applied-research candidate whose professional engineering experience, current
graduate research, and K--12 teaching preparation provide an equivalent path
into the work. I am available for remote collaboration during my final
semester and would be interested in a long-term research relationship as I
transition into doctoral study.

Thank you for considering this evidence-based equivalency case. I would value
the opportunity to discuss where my AI research, technical execution, and
education experience could contribute most effectively to DeepBrain's work.

Sincerely,\\
\textbf{Piter Z. Garcia Bautista}

\end{document}
